\documentclass[12pt,conference,onecolumn]{IEEEtran}

\usepackage[hidelinks]{hyperref}

\title{Teleoperated Robotic Arm Hand And Shoulder Exoskeleton}
\author{%
\IEEEauthorblockN{Andrew Dolgin}\IEEEauthorblockA{Science \& Engineering\\Manalapan High School\\Englishtown, NJ\\\href{mailto:427ssanrkitya@frhsd.com}{427ssankritya@frhsd.com}}\and
\IEEEauthorblockN{Julia Khabinskiy}\IEEEauthorblockA{Science \& Engineering\\Manalapan High School\\Englishtown, NJ\\\href{mailto:427mreznik@frhsd.com}{427mreznik@frhsd.com}}\and
\IEEEauthorblockN{Mason Levine}\IEEEauthorblockA{Science \& Engineering\\Manalapan High School\\Englishtown, NJ\\\href{mailto:427lbrunie@frhsd.com}{427lbrunie@frhsd.com}}
}

\date{January 28, 2027}

\newcommand{\keywords}{stroke, hemiparesis, servos, rehabilitation, TPU, PETG, flex sensors, microcontroller}

\usepackage{hyperref}
\makeatletter
\AtBeginDocument{
\hypersetup{%
pdftitle={\@title},
pdfauthor={Andrew Dolgin, Julia Khabinskiy, and Mason Levine},
pdfkeywords={\keywords}}}
\makeatother

\begin{document}
\maketitle 

\begin{abstract}
Every year, nearly 12 million people experience a stroke, and about 94 million are living as stroke survivors, making it the third leading cause of death and disability worldwide. Many survivors endure hemiparesis, which affects their arm and hand function. While consistent rehabilitation can improve recovery, existing devices are often inaccessible. High-end rehabilitation robots are too expensive, and affordable options typically lack effectiveness. This project aims to develop a portable, low-cost rehabilitation device that mirrors the motions of the patient’s unaffected side by using flex sensors, based on the principles of mirror therapy. We optimize an existing prototype glove by using flexible materials, like TPU, for a better fit and then design an articulated elbow and shoulder assembly. The device utilizes high-powered servos and a microcontroller to replicate natural shoulder and elbow movement, enhancing rehabilitation outside clinical settings.
\end{abstract}

\begin{IEEEkeywords}
\keywords
\end{IEEEkeywords}

\end{document}

